% ============================================================
% ch5_conclusion.tex
% ANSpect O-EMA · MS Thesis · UTD Electrical Engineering
% \input target — no \documentclass or \begin{document}
% STATUS: SKELETON — run 'expand CH-5.X' to flesh out
% v1.0  01-Sep-2026  Split from chapters.tex
% ============================================================

% ============================================================
% CHAPTER 5 — CONCLUSION AND FUTURE WORK
% ============================================================

\chapter{Conclusion and Future Work}
\label{ch:conclusion}

This thesis presented the design, implementation, and verification
of a J.83B Annex~B 256-QAM downstream transmit chain on the
Xilinx Zynq UltraScale+ \ac{rfsoc}~4x2.
The implemented system produced a correct \ac{rf} signal at
177\,MHz with verified spectral characteristics; consumer device
lock was not achieved, and the root cause remains an open problem.

% ------------------------------------------------------------
\section{Summary of Contributions}
\label{sec:contributions_summary}

\begin{tightemize}

  \item \textbf{Manual SystemVerilog \ac{rtl} implementation of
        the J.83B Annex~B 256-QAM downstream transmit chain.}
        Eight custom \ac{rtl} modules were designed and integrated
        as a ten-stage \ac{axi}-Stream pipeline on a 250\,MHz
        fabric clock, covering the full processing path from
        transport framing through \ac{srrc} pulse shaping.
        The manual coding approach was adopted after
        \ac{hdl} Coder was found to introduce fixed-point
        conversion overhead and integration constraints that made
        the generated \ac{rtl} unsuitable for direct deployment.

  \item \textbf{Multi-stage hardware verification methodology.}
        A stage-boundary golden vector framework was developed
        to enable bit-exact simulation and \ac{dma} loopback
        verification at each pipeline boundary.
        Unit simulation was applied to all eight custom modules;
        integration simulation was run across three incremental
        chain boundaries; and hardware golden vector comparison
        was performed at every \ac{dma}-accessible boundary.

  \item \textbf{\ac{rf} diagnostic methodology for \ac{rfsoc}
        transmit systems.}
        An \ac{rfdc} isolation test validated the analog signal
        path independently of the digital chain; the passthrough
        path achieved full consumer demodulator lock with
        signal strength, signal quality, and symbol quality all
        at 100\,\%.
        Thirteen hypotheses covering \ac{rfdc} configuration,
        \ac{iq} polarity, spectral inversion, constellation
        mapping, test vector content, and control word encoding
        were investigated and eliminated, localizing the consumer
        device lock failure to the protocol layer.

  \item \textbf{An integrated J.83B Annex~B 256-QAM downstream
        transmit system on \ac{rfsoc}.}
        The system was implemented as a PYNQ overlay on the
        RFSoC~4x2 and demonstrated continuous, gap-free
        \ac{axi}-Stream output at the \ac{rfdc} \ac{dac} input
        under sustained operation.
        A 131-second continuity measurement recorded
        4,219,220,880 samples with zero idle cycles across
        14 cyclic descriptor wrap transitions.

\end{tightemize}

% ------------------------------------------------------------
\section{Future Work}
\label{sec:future_work}

Three areas are identified for future investigation.

\textbf{Consumer device lock resolution.}
The consumer device lock failure was localized to the protocol
layer but the root cause was not identified within the scope
of this thesis.
Six technical hypotheses remain open: \ac{fec} frame
synchronization pattern acquisition requirements, constellation
point fixed-point precision compliance, \ac{srrc} interpolation
rate compliance, the clock tree and symbol rate relationship,
independent algorithmic verification of the constellation map,
and an independent audit of the \ac{rs} encoder, convolutional
interleaver, randomizer, and frame synchronization encoder
against the J.83B Annex~B specification.
Resolving these hypotheses requires a spec-compliant J.83B
software receiver capable of identifying the layer at which
the transmitted bitstream diverges from the standard, and a
cable signal analyzer with 256-QAM demodulation capability
for constellation visualization.

\textbf{Rx chain implementation.}
The scope of this thesis was restricted to the transmit
direction.
A complete J.83B Annex~B receiver chain on the same
\ac{rfsoc}~4x2 platform would enable closed-loop hardware
loopback testing, provide an independent method for
verifying transmit chain compliance, and complete the
bidirectional cable modem physical-layer prototype.

\textbf{O-EMA overlay activation.}
The ANSpect O-EMA \ac{phy} overlay is a proprietary
signal processing algorithm developed by ANSpect that
extends the J.83B Annex~B physical layer.
It was not activated within the scope of this thesis.
Integrating the O-EMA layer with the verified J.83B transmit
chain represents the next stage of the ANSpect platform
development.